%%%%%%%%%%%%%%%%%%%%%%%%%%%%%%%%%%%%%%%%%%%%%%%%%%%%%%%%%%%%%%%%%%%%%%%%%%%%%%%
% CHAPTER: The seven faces of r_{CY}(z) for Calabi-Yau chiral algebras
% Vol III specialization of the Vol I / Vol II seven-face programme.
% Mirrors the Vol I and Vol II seven-face master chapters.
%%%%%%%%%%%%%%%%%%%%%%%%%%%%%%%%%%%%%%%%%%%%%%%%%%%%%%%%%%%%%%%%%%%%%%%%%%%%%%%

%% Cross-volume label stubs. These labels are defined in Volume I and
%% referenced from this chapter. The phantomsection stubs prevent
%% undefined-reference errors at compile time; the canonical statements
%% live in Vol I.
\phantomsection\label{thm:gz26-commuting-differentials}%
\phantomsection\label{thm:gaudin-yangian-identification}%
\phantomsection\label{thm:yangian-sklyanin-quantization}%
\phantomsection\label{thm:dnp-bar-cobar-identification}%
\phantomsection\label{thm:kz-classical-quantum-bridge}%
\phantomsection\label{thm:shadow-depth-operator-order}%
\phantomsection\label{thm:vol1-seven-face-master}%
\phantomsection\label{thm:vol2-seven-face-master}%
\phantomsection\label{def:p-max}%
\phantomsection\label{def:k-max}%
\phantomsection\label{def:r-max}%

\chapter{The seven faces of $r_{CY}(z)$ for Calabi--Yau chiral algebras}
\label{ch:cy-holographic-datum-master}

\begin{quote}\itshape
For a Calabi--Yau chiral algebra arising from the CY-to-chiral functor,
the binary collision residue
$r_{CY}(z) = \Res^{\mathrm{coll}}_{0,2}(\Theta_{CY})$
has seven equivalent realizations specific to the CY setting: as the CY
bar-cobar twisting morphism, as a CoHA / perverse-coherent-sheaves shadow,
as the classical CY Poisson coisson, as the Maulik--Okounkov $R$-matrix
(for $K3 \times E$), as the affine super Yangian $Y(\widehat{\fgl}_1)$
$r$-matrix (for toric CY$_3$), as the elliptic Sklyanin bracket
(for toroidal CY), and as a Gaudin model arising from CY$_3$ collision
residues. The seven-face identification in
the CY setting specializes Vol~I's general theorems to the
CY$_3$ cases.
\end{quote}

%% =====================================================================
%% SECTION 1: The CY3 collision residue
%% =====================================================================

\section{The CY$_3$ collision residue}
\label{sec:cy3-collision-residue}

The collision residue is a single algebraic operation extracted from the
universal Maurer--Cartan element $\Theta_A$ of a chiral algebra. In the
CY setting, this operation acquires geometric meaning: the residue is
computed against the holomorphic CY volume form, and the resulting binary
operation is the classical limit of the quantum vertex chiral group
braiding.

\begin{definition}[Binary collision residue, CY version]
\label{def:cy-collision-residue}
Let $\cC$ be a smooth proper Calabi--Yau category of dimension~$d$ with
holomorphic volume form $\Omega_\cC$. Let $A_\cC$ denote the chiral
algebra produced by the CY-to-chiral functor (Theorem~CY-A$_2$ at
$d = 2$; Theorem~CY-A$_3$ at $d = 3$, Theorem~\ref{thm:cy-to-chiral-d3}). Let
$\Theta_{A_\cC} \in \mathrm{MC}(\fg^{\mathrm{mod}}_{A_\cC})$
denote its universal Maurer--Cartan element (cf.\ Vol~I,
Theorem~\ref{thm:bar-cobar-adjunction}). The \emph{binary CY collision
residue} is
\[
 r_{CY}(z) \;:=\; \Res^{\mathrm{coll}}_{0,2}\bigl(\Theta_{A_\cC}\bigr)
 \;\in\; A_\cC \otimes A_\cC \,[\![z^{-1}]\!],
\]
the genus-$0$, degree-$2$ residue along the collision divisor of
$\overline{C}_2(\Spec\C)$, paired against $\Omega_\cC$ in the CY
direction.
\end{definition}

\begin{remark}[The CY direction is a separate slot]
\label{rem:cy-direction-slot}
The collision residue extracts the $z^{-1}$ part of $\Theta$ along the
chiral curve direction. In the CY setting, $\Theta$ has \emph{two}
gradings: degree in the chiral direction (controlled by the curve $X$)
and Hodge weight in the CY direction (controlled by~$\cO_\cC$). The CY
volume form $\Omega_\cC$ kills the Hodge weight by integration; what
remains is a chiral binary operation valued in $A_\cC \otimes A_\cC$.
This is the precise sense in which $r_{CY}$ couples the two directions.
\end{remark}

\begin{remark}[Three invariants must not be conflated]
\label{rem:cy-three-invariants}
For any chiral algebra $A$, the binary residue carries three
\emph{distinct} numerical invariants which the CY setting will
illuminate:
\begin{enumerate}[label=(\roman*)]
\item $p_{\max}(A)$, the maximal pole order of generator-on-generator
OPEs (cf.\ Vol~I, Definition~\ref{def:p-max});
\item $k_{\max}(A) = p_{\max}(A) - 1$, the collision depth of $r_{CY}$
itself, after $d\log$-absorption (cf.\ Vol~I,
Definition~\ref{def:k-max});
\item $r_{\max}(A)$, the shadow depth, the degree at which the
obstruction tower $\Theta_A^{\leq r}$ terminates (Vol~I,
Definition~\ref{def:r-max}).
\end{enumerate}
The $\beta\gamma$ system is the archetypal witness of their independence:
$p_{\max}(\beta\gamma) = 1$, $k_{\max}(\beta\gamma) = 0$,
$r_{\max}(\beta\gamma) = 4$ (class~C). Conflating these produces
incorrect classifications, and the seven-face programme will distinguish
them at every face.
\end{remark}

%% =====================================================================
%% SECTION 2: Face 1 -- CY bar-cobar twisting morphism
%% =====================================================================

\section{Face 1: the CY bar-cobar twisting morphism}
\label{sec:face1-bar-cobar}

The first face is the home framework of this volume. The CY-to-chiral
functor sends a Calabi--Yau category $\cC$ to a chiral algebra $A_\cC$,
and the bar construction sends $A_\cC$ to a factorization coalgebra
$\barB(A_\cC)$ on $\Ran(X)$. The collision residue is the binary
component of the universal twisting morphism between $\barB(A_\cC)$ and
$A_\cC$.

\begin{theorem}[Face 1: CY bar-cobar realization, $d = 2$]
\label{thm:face1-cy-bar-cobar}
\ClaimStatusProvedHere
Let $A_\cC$ be the chiral algebra of a CY$_2$ category $\cC$ produced by
the CY-to-chiral functor $\Phi$ (Theorem~CY-A$_2$). Then:
\begin{enumerate}[label=(\roman*)]
\item The convolution dg Lie algebra $\fg^{\mathrm{mod}}_{A_\cC}$
identifies with the homotopy Lie algebra of twisting morphisms
$\mathrm{hom}_\alpha(\barB(A_\cC), A_\cC)$ between bar coalgebra and
chiral algebra (cf.\ Vol~I,
Theorem~\ref{thm:bar-cobar-adjunction});
\item The universal MC element $\Theta_{A_\cC}$ is the universal
twisting morphism;
\item The binary collision residue $r_{CY}(z) =
\Res^{\mathrm{coll}}_{0,2}(\Theta_{A_\cC})$ is the leading
degree-$2$ component of this twisting morphism, evaluated in the
CY direction against $\Omega_\cC$.
\end{enumerate}
\end{theorem}

\begin{proof}[Proof sketch]
For~(i), the cyclic bar complex $\mathrm{CC}_\bullet(\cC)$ identifies
with $\barB(A_\cC)$ as a factorization coalgebra
(Proposition~\ref{prop:bar-dictionary}). The convolution identification is
the application of Vol~I's master twisting-morphism theorem to this
coalgebra-algebra pair. For~(ii), $\Theta_{A_\cC}$ satisfies the MC
equation $D\Theta + \tfrac{1}{2}[\Theta,\Theta] = 0$ at all levels, with
the bar-intrinsic construction $\Theta = D - d_0$. For~(iii), the
collision residue is computed from $\Theta$ by restricting to genus~$0$,
degree~$2$, and pairing against the holomorphic volume form.
\end{proof}

\begin{theorem}[Face 1: CY bar-cobar realization, $d = 3$; \ClaimStatusProvedHere]
\label{conj:face1-cy-bar-cobar-d3}
The statement of Theorem~\ref{thm:face1-cy-bar-cobar} extends to
$d = 3$ by Theorem~CY-A$_3$ (Theorem~\ref{thm:cy-to-chiral-d3}).
The CY$_3$-to-chiral functor $\Phi_3$ is defined on smooth proper CY$_3$ categories,
and for every such $\cC$ the convolution
identification, the universal MC element, and the binary collision
residue extraction of Theorem~\ref{thm:face1-cy-bar-cobar}(i)--(iii)
hold for $A_\cC := \Phi_3(\cC)$.
\end{theorem}

\begin{remark}[Bar-cobar inversion is not the seven-face master move]
\label{rem:no-cobar-bulk-confusion}
The cobar functor $\Omega$ inverts the bar complex: $\Omega(\barB(A))
\simeq A$ (Vol~I, Theorem~B). This is a round-trip that recovers $A$
itself. The seven-face programme is not about recovering $A$ from $\barB(A)$;
it is about realizing a single algebraic object, the binary residue
$r_{CY}(z)$, in seven distinct geometric languages. The closed-string
or bulk side is the chiral derived center $\cZ^{\mathrm{der}}_{\mathrm{ch}}(A_\cC)$
(Vol~I, Theorem~H), not the cobar; the seven faces all live on the
boundary~$A_\cC$ side.
\end{remark}

\begin{remark}[BPS shadow depth and the CY holographic datum]
\label{rem:bps-shadow-depth-cy}
\ClaimStatusProvedElsewhere
The holographic modular Koszul datum in the CY setting assembles the six
pieces of the seven-face programme into a single package:
\[
 H_{CY}(T) \;=\; \bigl(A_{CY},\; A_{CY}^{!},\; \cC_{CY},\; r_{CY}(z),\;
 \Theta_{CY},\; \nabla^{\mathrm{hol}}_{CY}\bigr),
\]
where $A_{CY}$ is the CY-to-chiral image $A_\cC = \Phi(\cC)$, $A_{CY}^{!}$
is its Verdier/Koszul dual boundary algebra, $\cC_{CY}$ is the line
category of the HT theory, $r_{CY}(z)$ is the binary CY collision residue
of Definition~\ref{def:cy-collision-residue}, $\Theta_{CY}$ is the
universal Maurer--Cartan element (Theorem~\ref{thm:face1-cy-bar-cobar}),
and $\nabla^{\mathrm{hol}}_{CY}$ is the holomorphic Knizhnik--Zamolodchikov
connection on the chiral curve direction. This is the CY specialization
of the Vol~I datum $H(T) = (A, A^{!}, C, r(z), \Theta_A, \nabla^{\mathrm{hol}})$;
the new feature is the presence of the CY direction $\Omega_\cC$
(Remark~\ref{rem:cy-direction-slot}).

Dimofte's framework (PIRSA 25110067) attaches to $H_{CY}(T)$ a single
numerical invariant, the \emph{BPS complexity} of the boundary Hilbert
space, measured by the $K$-matrix $K_{A_{CY}}(z)$ of
equation~\eqref{eq:dimofte-k-matrix-cy}
(cf.\ Chapter~\ref{ch:quantum-chiral-algebras},
Remark~\ref{rem:dimofte-k-matrix-cy}, and \emph{Vol~II}, remark
\texttt{rem:dimofte-k-matrix} in
\texttt{ordered\_associative\_chiral\_kd\_core.tex}). The boundary
Hilbert space of the HT theory decomposes as
\[
 \cH_{\partial}(T) \;\simeq\; \cH_{\partial}^{\mathrm{free}} \otimes
 \cH_{\partial}^{\mathrm{BPS}},
\]
where $\cH_{\partial}^{\mathrm{free}}$ is the $\mathfrak{u}(1)^r$ free-boson
factor (the class $\mathbf{G}$ summand) and $\cH_{\partial}^{\mathrm{BPS}}$
records the BPS corrections. Under the CY-to-chiral functor, this
decomposition is the Vol~I shadow partition $\mathbf{G} +
(\mathbf{L}/\mathbf{C}/\mathbf{M})$: the free-boson summand is exactly
class $\mathbf{G}$, and the BPS summand carries the higher shadow depth.

\paragraph{CY$_2$ (K3, proved).} For a CY$_2$ category $\cC$ admitting a
chiral algebraization via $\Phi$, the BPS Hilbert space is determined by
the associated Borcherds--Kac--Moody superalgebra
$\mathfrak{g}_{\mathrm{BKM}}(\cC)$, and the $K$-matrix encodes its
denominator identity. Concretely, for $\cC = D^b(\mathrm{Coh}(K3))$ the
BKM denominator formula (Gritsenko--Nikulin, Borcherds) identifies the
character of $\cH_{\partial}^{\mathrm{BPS}}$ with the Fourier coefficients
of $1/\Phi_{10}$, and the $K$-matrix modes $\{\alpha_n\}$ read off the
root multiplicities recorded by these coefficients.

\paragraph{Igusa cusp form verification.} The Igusa cusp form $\Phi_{10}$
on $\mathfrak{H}_2$ has weight $10$ (Gritsenko--Nikulin). The
Borcherds--Kac--Moody modular characteristic of $K3 \times E$ is
$\kappa_{\mathrm{BKM}} = 5$. The
identification
\[
 \mathrm{wt}(\Phi_{10}) \;=\; 2\,\kappa_{\mathrm{BKM}}
 \;=\; 10
\]
holds: the factor of $2$ is the Serre duality on $K3$, which has Serre
degree $d = 2$, so the BKM characteristic is one half of the cusp-form
weight. This matches entry $\kappa_{\mathrm{BKM}}(K3 \times E) = 5$ in
the Vol~III $\kappa_{\mathrm{BKM}}$-spectrum table, and agrees with the identity
$\mathrm{wt}(\Phi_{10}) = 10$ recorded in
Chapter~\ref{ch:k3-times-e}. AP-CY8 forbids calling this an identity at
the theorem level: it is an observation matching the Borcherds lift
weight, not a computation from the Vol~I definition of
$\kappa_{\mathrm{ch}}$.

\paragraph{CY$_3$ (CY-A$_3$ proved).} For a CY$_3$ category the
BPS Hilbert space should be the module over the Donaldson--Thomas Hall
algebra (Kontsevich--Soibelman, Davison--Meinhardt), and the $K$-matrix
should be the generating function of the cohomological Hall algebra
$\cH(Q, W)$ attached to a quiver-with-potential presentation. For toric
CY$_3$ without compact $4$-cycles (Chapter~\ref{ch:toric-coha}) this is
explicit: Schiffmann--Vasserot identify $\cH(\C^3)$ with the positive
half of the affine Yangian $Y^{+}(\widehat{\mathfrak{gl}}_1)$, and the
$K$-matrix is its character $M(q)^2 \cdot P(q)$
(Theorem~\ref{thm:drinfeld-center-coha}). All statements in this paragraph
whose proof chain passes through $A_\cC$ at $d = 3$ follow from
Theorem~CY-A$_3$ (Theorem~\ref{thm:cy-to-chiral-d3}, proved).

\paragraph{Shadow--BPS dictionary.} The four-class partition of CY
boundary algebras is:
\begin{center}
\renewcommand{\arraystretch}{1.2}
\begin{tabular}{llll}
\toprule
Shadow class & BPS complexity & $K_{A_{CY}}(z)$ & Example \\
\midrule
$\mathbf{G}$ & 0 BPS corrections & $1$ & trivial CY fiber, $\mathfrak{u}(1)^r$ \\
$\mathbf{L}$ & 1 BPS level & polynomial, simple pole & rank-one CY Heisenberg deformation \\
$\mathbf{C}$ & 2 BPS levels & polynomial, double pole & toric CY$_3$ ($\C^3$, conifold) \\
$\mathbf{M}$ & infinite BPS tower & infinite expansion & $K3$, $K3 \times E$ \\
\bottomrule
\end{tabular}
\end{center}

The shadow-depth classification of the CY boundary algebra is therefore
equivalent to the BPS complexity of the HT boundary Hilbert space under
the CY-to-chiral functor. The $K$-matrix modifies the coproduct, not the
product, so the level-$k$ vanishing check does not apply directly;
the companion $r$-matrix check is the level check on the affine
$r$-matrix at $k_{\mathrm{ch}} = 0$, which vanishes as required and
returns class $\mathbf{G}$. Primary references for BPS Hall algebras and
the denominator identity: Dimofte PIRSA 25110067, Kontsevich--Soibelman,
Davison--Meinhardt, Gritsenko--Nikulin. Vol~II cross-references:
\texttt{rem:dimofte-k-matrix} in
\texttt{ordered\_associative\_chiral\_kd\_core.tex} and the BPS/shadow
discussion in \texttt{ht\_bulk\_boundary\_line\_core.tex}.
\end{remark}

%% =====================================================================
%% SECTION 3: Face 2 -- CoHA / perverse coherent sheaves
%% =====================================================================

\section{Face 2: CoHA and perverse coherent sheaves}
\label{sec:face2-coha}

The second face is enumerative geometry: the binary residue is the leading
non-trivial structure constant of the cohomological Hall algebra of $\cC$,
or equivalently, the Yoneda product on perverse coherent sheaves. This
face is the natural home of Donaldson--Thomas invariants.

\begin{conjecture}[Face 2: CoHA realization of $r_{CY}$]
\label{conj:face2-coha-realization}
Let $\cC$ be a smooth proper CY$_3$ category with quiver-with-potential
presentation $(Q, W)$ at a stable point of the moduli of stability
conditions. Let $\cH(Q,W)$ denote the critical cohomological Hall
algebra (Definition~\ref{def:critical-coha}). Then the binary collision
residue of $A_\cC$ admits a CoHA realization
\[
 r_{CY}^{\mathrm{CoHA}}(z) \;=\; \mu_{\cH(Q,W)}^{(2)} \otimes z^{-1}
 \;\in\; \cH(Q,W) \otimes \cH(Q,W) \,[\![z^{-1}]\!],
\]
where $\mu^{(2)}$ is the binary CoHA product, paired with the
$d\log$-propagator on $\overline{C}_2$.
\ClaimStatusConjectured
\end{conjecture}

\begin{remark}[Status of Conjecture~\ref{conj:face2-coha-realization}]
\label{rem:face2-status}
The CoHA $\cH(Q,W)$ exists unconditionally (Kontsevich--Soibelman,
Davison--Meinhardt). The chiral algebra $A_\cC$ exists for $d = 2$ and
is conjectural for $d = 3$. The identification with $r_{CY}$ requires
the CY-to-chiral functor to be defined and to commute with critical
cohomology, which is the content of Conjecture~CY-C and is currently
verified for the toric CY$_3$ examples
(Theorems~\ref{thm:sv-c3}, \ref{thm:rsyz}). In those cases the
identification reduces to the Schiffmann--Vasserot and
Rapcak--Soibelman--Yang--Zhao theorems and is ProvedElsewhere with
attribution to those two papers.
\end{remark}

\begin{remark}[Perverse coherent sheaves]
\label{rem:perverse-coherent-sheaves}
For a CY$_3$ category arising as $D^b(\Coh(X))$ for $X$ a smooth CY$_3$
variety, the CoHA description above is equivalent to a description in
terms of perverse coherent sheaves on $X$ (Bridgeland, Bayer--Macri).
The Yoneda product on $\Ext^{\bullet}$ of perverse coherent sheaves is
the binary CoHA product, and the binary collision residue extracts its
leading $\Ext^1$ component. This is the geometric face: the residue
$r_{CY}$ \emph{is} a piece of the deformation theory of the CY$_3$.
\end{remark}

\begin{remark}[CoHA is not the full quantum vertex chiral group]
\label{rem:coha-not-qvcg}
The critical CoHA is an associative algebra: the positive (or
$E_1$-ordered) half of the conjectural quantum vertex chiral group
$G(X)$. The full $G(X)$ would be braided $E_2$, with the braiding
supplied by the $R$-matrix of the affine super Yangian
(Section~\ref{sec:face5-yangian}). The CoHA face captures the
\emph{shadow} of $r_{CY}$ on the ordered side; the braided side is
faces~4 and~5.
\end{remark}

%% =====================================================================
%% SECTION 4: Face 3 -- Classical CY Poisson coisson
%% =====================================================================

\section{Face 3: classical CY Poisson coisson \textup{(}genus~$0$ only\textup{)}}
\label{sec:face3-coisson}

The third face is the classical limit. In the CY setting, the binary
residue $r_{CY}$ has a classical avatar: a coisson bracket on the
sheafification of $A_\cC$, which in turn is a chiral version of the
holomorphic Poisson bracket arising from the CY $(2{,}0)$-form (in
$d = 2$) or the holomorphic top form via contraction (in $d = 3$).

\begin{theorem}[Face 3: classical Poisson coisson, $d = 2$ \textup{(}genus~$0$ only\textup{)}]
\label{thm:face3-classical-coisson}
\ClaimStatusProvedHere
Let $\cC$ be a CY$_2$ category and let $A_\cC$ be its chiral
algebraization (Theorem~CY-A$_2$). Let $\hbar$ denote a quantization
parameter. Then:
\begin{enumerate}[label=(\roman*)]
\item The classical limit $A_\cC^{\mathrm{cl}} := \lim_{\hbar \to 0}
A_\cC$ is a Poisson vertex algebra (PVA) with bracket
$\{\cdot,\cdot\}_\lambda$ valued in $A_\cC^{\mathrm{cl}} \otimes
\C[\lambda]$;
\item The leading coefficient $\{\cdot,\cdot\}_\lambda^{(0)}$ of the
PVA bracket is a chiral Poisson bracket whose classical $r$-matrix
is exactly $r_{CY}(z)$ at $\hbar = 0$;
\item This identifies $r_{CY}^{\mathrm{cl}}$ as the coisson bracket
constructed from the CY volume form $\Omega_\cC$ via the Poisson
structure on $\cZ^{\mathrm{der}}_{\mathrm{ch}}(A_\cC^{\mathrm{cl}})$
(cf.\ Vol~I, Theorem~\ref{thm:kz-classical-quantum-bridge}).
\end{enumerate}
\end{theorem}

\begin{proof}[Proof sketch]
For a CY$_2$ category, $\Phi$ is the proved CY-to-chiral functor
(Chapter~\ref{ch:cy-to-chiral}). The PVA structure on $A_\cC^{\mathrm{cl}}$
is the classical limit of the chiral OPE. The leading $\lambda$
coefficient of the $\lambda$-bracket is the contraction of two factors
along the binary collision divisor; this is exactly
$\Res^{\mathrm{coll}}_{0,2}$ computed in the classical limit, hence
$r_{CY}^{\mathrm{cl}}$. The identification with the coisson bracket
arising from $\Omega_\cC$ is the coisson--PVA dictionary
(cf.\ Vol~I, Theorem~\ref{thm:kz-classical-quantum-bridge}, applied to
the specialization $\cC \mapsto A_\cC$).
\end{proof}

\begin{conjecture}[Face 3: classical Poisson coisson, $d = 3$ \textup{(}genus~$0$ only\textup{)}]
\label{conj:face3-classical-coisson-d3}
The statement of Theorem~\ref{thm:face3-classical-coisson} extends to
$d = 3$, by Theorem~CY-A$_3$ (Theorem~\ref{thm:cy-to-chiral-d3}): if $A_\cC$ is the
CY$_3$-chiral algebraization of $\cC$, then items (i)--(iii) of
Theorem~\ref{thm:face3-classical-coisson} hold for $A_\cC^{\mathrm{cl}}$
with the CY$_3$ volume form $\Omega_\cC \in H^0(\cC, \omega_\cC)$
supplying the coisson pairing.
\ClaimStatusConjectured
\end{conjecture}

\begin{remark}[Convention: $\lambda$-brackets vs OPE modes]
\label{rem:lambda-vs-mode}
Vol~I uses OPE modes $a_{(n)}b$. Vol~II uses $\lambda$-brackets
$\{a_\lambda b\} = \sum_n \lambda^{(n)} a_{(n)}b$ with the divided
power $\lambda^{(n)} = \lambda^n/n!$. Vol~III uses motivic and categorical
language. When transcribing $r_{CY}$ between volumes, the $1/n!$ factor
must be tracked: the $\lambda^3$ coefficient of a $\lambda$-bracket is
$a_{(3)}b/3! = a_{(3)}b/6$, not $a_{(3)}b$. This convention bridge is
the source of several historical errors and is discussed in the cross-volume convention bridge (Vol~I, concordance).
\end{remark}

%% =====================================================================
%% SECTION 5: Face 4 -- Maulik-Okounkov R-matrix
%% =====================================================================

\section{Face 4: the Maulik--Okounkov $R$-matrix (for $K3 \times E$)}
\label{sec:face4-mo-rmatrix}

The fourth face is geometric representation theory: the binary residue is
the classical limit of the Maulik--Okounkov stable-envelope $R$-matrix
acting on the equivariant cohomology of a Nakajima quiver variety. For
the CY$_3$ specialization $K3 \times E$, this is the most concrete
realization, and it produces $r_{CY}$ explicitly in terms of theta
functions on the elliptic factor.

The Maulik--Okounkov stable envelope and its $R$-matrix
$R_{\mathrm{MO}}^{K3 \times E}(z)$ exist unconditionally (Maulik--Okounkov
2019); this is the ProvedElsewhere component. The genuinely new content
of face~4 is the identification of the classical limit with the binary
CY collision residue.

\begin{conjecture}[Face 4: MO realization for $K3 \times E$]
\label{conj:face4-mo-k3e}
For the CY$_3$ specialization $X = K3 \times E$, the binary CY
collision residue identifies with the classical limit of the
Maulik--Okounkov $R$-matrix:
\[
 r_{CY}^{K3 \times E}(z) \;=\;
 \lim_{\hbar \to 0} \, \hbar^{-1}\bigl(R_{\mathrm{MO}}^{K3 \times E}(z)
 - \id\bigr).
\]
The right-hand side is the rational/elliptic stable envelope of
Maulik--Okounkov, evaluated on the Hilbert scheme of points on $K3$
along the elliptic direction $E$. The conjecture is equivalent to
Conjecture~CY-C for $K3 \times E$.
\ClaimStatusConjectured
\end{conjecture}

\begin{remark}[The CY involution $g(z)g(-z) = 1$]
\label{rem:cy-involution-rmatrix}
For $K3 \times E$, the MO $R$-matrix is governed by a single Yangian
generator $g(z)$ satisfying $g(z) g(-z) = 1$
(Theorem~\ref{thm:k3e-mo-rmatrix}). This CY involution is the
$\mathbb{S}^2$-framing of the CY-to-chiral functor seen at the level of
the binary residue: the $z \mapsto -z$ involution is the chiral analogue
of the CY Serre functor $\mathbb{S} \simeq [d]$. In particular the CY
involution forces the residue to be of \emph{odd} pole order, in
agreement (the $d\log$-absorption sends even-order poles to
odd-order poles in the residue).
\end{remark}

\begin{remark}[Self-dual Gepner trivialization]
\label{rem:gepner-trivialization}
At the Gepner point of $K3$, where the chiral algebra has $\cN = (4,4)$
enhancement, the MO $R$-matrix trivializes
(Conjecture~\ref{conj:k3e-self-dual-limit}), so $r_{CY}$ vanishes.
This is consistent with the seven-face master statement: at the self-dual
point of any face, all seven realizations vanish simultaneously. For
$K3 \times E$ this is a statement about the vanishing of the collision
residue in the enhanced-symmetry limit, not an identification of the
global modular characteristic with zero.
\end{remark}

%% =====================================================================
%% SECTION 6: Face 5 -- Affine super Yangian Y(gl_1) for toric CY3
%% =====================================================================

\section{Face 5: the affine super Yangian $Y(\widehat{\fgl}_1)$ for toric CY$_3$}
\label{sec:face5-yangian}

The fifth face is the explicit toric face: when the CY$_3$ is toric and
admits a quiver-with-potential presentation, the binary residue is the
classical $r$-matrix of the affine super Yangian
$Y(\widehat{\fgl}_1)$ (for $\C^3$, the Jordan quiver) or more generally
$Y(\widehat{\fg}_{Q_X})$ (for the toric quiver of~$X$).

\begin{theorem}[Face 5 for toric CY$_3$ via RSYZ]
\label{thm:face5-toric-yangian-rsyz}
\ClaimStatusProvedElsewhere
Let $X$ be a toric CY$_3$ without compact $4$-cycles, with toric quiver
$(Q_X, W_X)$ and affine super Yangian $Y(\widehat{\fg}_{Q_X})$
(Theorem~\ref{thm:rsyz}). Then the binary CY collision residue
identifies with the classical $r$-matrix:
\[
 r_{CY}^{X}(z) \;=\; r_{Y(\widehat{\fg}_{Q_X})}^{\mathrm{cl}}(z)
 \;=\; \frac{\Psi\,\Omega_{Y(\widehat{\fg}_{Q_X})}}{z}
 \;+\; (\text{higher poles dictated by } p_{\max}(W_X)),
\]
where $\Omega_{Y(\widehat{\fg}_{Q_X})}$ is the Casimir tensor of the
affine super Yangian. At the Jordan-quiver specialization $X = \C^3$,
the identification reduces to the $r$-matrix of the affine Yangian
$Y(\widehat{\fgl}_1) \cong \cW_{1 + \infty}$ at the self-dual level.
\ClaimStatusProvedElsewhere
\end{theorem}

\begin{remark}[Status partition watch]
\label{rem:face5-ap60}
Theorem~\ref{thm:face5-toric-yangian-rsyz} is a restatement of the
Rapcak--Soibelman--Yang--Zhao theorem (Theorem~\ref{thm:rsyz}) in the
collision-residue language of the seven-face programme. The status
partition principle (AP60) forbids tagging this theorem ProvedHere: the classical content is due to
Schiffmann--Vasserot (for $\C^3$, Theorem~\ref{thm:sv-c3}) and
Rapcak--Soibelman--Yang--Zhao (for general toric without compact
$4$-cycles, Theorem~\ref{thm:rsyz}). The genuinely new content of
face~5 is the translation of their $r$-matrix into the language of
$\Res^{\mathrm{coll}}_{0,2}(\Theta_{A_X})$, which is recorded as the
scholium below.
\end{remark}

\begin{proposition}[Collision-residue translation of face 5, $\C^3$ case]
\label{prop:face5-c3-translation}
For $X = \C^3$, the Schiffmann--Vasserot Casimir tensor
$\Omega_{Y(\widehat{\fgl}_1)}$ coincides with
$\Res_{z=0}\bigl[r_{CY}^{\C^3}(z)\bigr]$, where $r_{CY}^{\C^3}$ is the
binary CY collision residue of the $\C^3$ chiral algebra.
\ClaimStatusProvedHere
\end{proposition}

\begin{proof}[Proof sketch]
The Jordan quiver superpotential $W = \tr(XYZ - XZY)$ has
$p_{\max}(W) = 1$ (each cubic monomial contributes a single contraction
in the binary OPE), so $k_{\max} = 0$ and $r_{CY}^{\C^3}(z)$ has a
\emph{single} pole at $z = 0$. The residue at this pole is the
Schiffmann--Vasserot Casimir
$\Omega_{Y(\widehat{\fgl}_1)} \in Y^+(\widehat{\fgl}_1)^{\otimes 2}$
(Theorem~\ref{thm:sv-c3}). By the universality of $\Theta$ in MC moduli,
this is exactly $r_{CY}^{\C^3}$. The matching of higher poles with the
toric quiver superpotential follows from the RSYZ theorem for toric
quivers without compact $4$-cycles.
\end{proof}

\begin{remark}[Three invariants for toric CY$_3$]
\label{rem:toric-three-invariants}
For the standard toric CY$_3$ family the three invariants are:
\begin{center}
\small
\begin{tabular}{lcccc}
\toprule
Geometry & Quiver & $p_{\max}$ & $k_{\max}$ & $r_{\max}$ \\
\midrule
$\C^3$ & Jordan & $1$ & $0$ & $2$ (class~G) \\
Resolved conifold & Klebanov--Witten & $1$ & $0$ & $2$ (class~G) \\
Local $\bP^2$ & McKay $\Z_3$ & $2$ & $1$ & $\infty$ (class~M) \\
Local $\bP^1\times\bP^1$ & Beilinson & $2$ & $1$ & $\infty$/$2$ \\
\bottomrule
\end{tabular}
\end{center}
The pole order $p_{\max}$, the collision depth $k_{\max} = p_{\max} - 1$,
and the shadow depth $r_{\max}$ are three independent invariants. In the
toric family above, $p_{\max}$ is small (at most $2$) but $r_{\max}$ can
be infinite (local $\bP^2$, McKay $\Z_3$). This separates the pole-order
face (face~5, controlled by $p_{\max}$) from the shadow tower face
(face~1, controlled by $r_{\max}$). AP59 forbids identifying these.
\end{remark}

\subsection{Gaudin Hamiltonians for toric CY$_3$ from collision residues}
\label{subsec:toric-gaudin-from-collisions}

The collision residue at the Jordan quiver $(\C^3)$ is the
$Y(\widehat{\fgl}_1)$ classical $r$-matrix; this allows one to write
down a chiral Gaudin model whose Hamiltonians are extracted from CY$_3$
collision data, providing the first instance of face~7 inside face~5.
The full development of face~7 occupies
Section~\ref{sec:face7-cy3-gaudin}.

\begin{construction}[Gaudin Hamiltonians from $r_{CY}^{\C^3}$]
\label{constr:c3-gaudin-from-residue}
Fix marked points $z_1, \ldots, z_n \in \C$ (the chiral curve, with $n$
copies of the Fock representation of $Y(\widehat{\fgl}_1)$). Define
\[
 H_i^{\C^3} \;:=\; \sum_{j \neq i} \frac{k\,\Omega_{ij}^{Y(\widehat{\fgl}_1)}}{z_i - z_j}
 \;\in\; Y(\widehat{\fgl}_1)^{\otimes n},
\]
where $\Omega_{ij}^{Y(\widehat{\fgl}_1)}$ is the $Y(\widehat{\fgl}_1)$
Casimir $\Res_{z=0} r_{CY}^{\C^3}(z)$, acting on the $i$-th and $j$-th
tensor factors. The operators $\{H_i^{\C^3}\}$ commute pairwise as a
direct consequence of the classical Yang--Baxter equation for
$r_{CY}^{\C^3}$ (cf.\ Vol~I,
Theorem~\ref{thm:gaudin-yangian-identification}).
\end{construction}

\begin{remark}[Wall-crossing as gauge equivalence]
\label{rem:c3-wall-crossing-mc-gauge}
Across walls in the Bridgeland stability manifold, the toric Yangian
$r$-matrix transforms by a Kontsevich--Soibelman wall-crossing
automorphism. In the present language, this is exactly an MC gauge
equivalence on $\Theta_{A_X}$ (Theorem~\ref{thm:wall-crossing-mc}):
$r_{CY}^{X,\zeta'}(z) = e^{\ad_\alpha} r_{CY}^{X,\zeta}(z)$ for the
unique gauge parameter $\alpha \in \fn_W$ supported on the wall divisor.
The Gaudin Hamiltonians of
Construction~\ref{constr:c3-gaudin-from-residue} are gauge-equivalent
across walls; their joint spectrum is chamber-independent.
\end{remark}

%% =====================================================================
%% SECTION 7: Face 6 -- Elliptic Sklyanin bracket
%% =====================================================================

\section{Face 6: elliptic Sklyanin bracket (for toroidal CY)}
\label{sec:face6-sklyanin}

The sixth face is the elliptic deformation: when the chiral curve is
itself an elliptic curve $E_\tau$ and the CY data are toroidal, the
binary residue acquires an elliptic structure governed by theta functions.
The classical limit of this $r$-matrix is the elliptic Sklyanin bracket,
and the quantization is Felder's elliptic quantum group.

The elliptic $r$-matrix and its classical Sklyanin bracket are due to
Sklyanin (1983) and Belavin--Drinfeld (1984), with the quantization to
Felder's elliptic quantum group (Definition~\ref{def:elliptic-quantum})
classical. These components are ProvedElsewhere with explicit
attribution. The genuinely new content is the identification of the
elliptic $r$-matrix with the binary CY collision residue
$\Res^{\mathrm{coll}}_{0,2}(\Theta_{A_\cC})$ via the chiral propagator
$d \log \theta_1(z\,|\,\tau)$.

\begin{theorem}[Face 6 for $\fsl_2$ Heisenberg on $E_\tau$]
\label{thm:face6-elliptic-sklyanin}
\ClaimStatusProvedHere
Let $E_\tau$ be an elliptic curve and let $A = H^{\fsl_2}_{E_\tau}$ be
the $\fsl_2$ Heisenberg chiral algebra realized on $E_\tau$. Then the
binary CY collision residue admits the elliptic representative
\[
 r_{CY}^{\tau}(z) \;=\;
 \frac{k\,\theta'_1(0\,|\,\tau)}{\theta_1(z\,|\,\tau)}\,\Omega_{\fsl_2},
\]
where $\theta_1$ is the odd Jacobi theta function and
$\Omega_{\fsl_2}$ is the Casimir tensor of~$\fsl_2$
(cf.\ Section~\ref{sec:elliptic-r-theta}). The classical limit of
$r_{CY}^{\tau}$ defines the Sklyanin bracket
$\{f,g\}_\tau = \langle r_{CY}^{\tau}, [\nabla f, \nabla g]\rangle$ on
$E_{q,p}(\fsl_2)^*$, and the quantization $\hbar \to q$ recovers
Felder's elliptic quantum group $E_{q,p}(\fsl_2)$ with dynamical
parameter $\lambda$ identified as the $B$-cycle monodromy of the
elliptic propagator.
\end{theorem}

\begin{conjecture}[Face 6 for general toroidal CY]
\label{conj:face6-general-toroidal}
The statement of Theorem~\ref{thm:face6-elliptic-sklyanin} extends to
general toroidal CY realizations $X$ on $E_\tau$: the binary residue
$r_{CY}^{\tau,X}(z)$ is the elliptic $r$-matrix of
$Y(\widehat{\fg}_X)$ acquiring dynamical theta-function corrections in
$\lambda$, and its classical limit is the corresponding Sklyanin
bracket (cf.\ Vol~I,
Theorem~\ref{thm:yangian-sklyanin-quantization}).
\ClaimStatusConjectured
\end{conjecture}

\begin{remark}[status partition]
\label{rem:face6-ap60}
The Sklyanin bracket and Belavin--Drinfeld classification are classical
results (1983). The genuinely new content of
Theorem~\ref{thm:face6-elliptic-sklyanin} is the identification of the
elliptic $r$-matrix with the binary CY collision residue
$\Res^{\mathrm{coll}}_{0,2}(\Theta_{A_\cC})$ via the chiral propagator
$d \log \theta_1(z\,|\,\tau)$. AP60 forbids tagging the entire theorem
ProvedHere; the classical components are ProvedElsewhere with explicit
attribution to Sklyanin (1983) and Belavin--Drinfeld (1984).
\end{remark}

\subsection{The three-parameter $\hbar$ identification}
\label{subsec:three-hbar}

A subtle but central observation is that three different $\hbar$
parameters appearing in the literature on elliptic quantum groups are
the same parameter under the seven-face dictionary.

\begin{construction}[Three-parameter $\hbar$ identification]
\label{constr:three-parameter-hbar}
The following three quantization parameters coincide:
\begin{enumerate}[label=(\roman*)]
\item the boundary coupling $\hbar^{\mathrm{KZ}}$ in the Knizhnik--Zamolodchikov
equation (Knizhnik--Zamolodchikov 1984; cf.\ Vol~I,
Theorem~\ref{thm:kz-classical-quantum-bridge});
\item the loop parameter $\hbar^{\mathrm{DNP}}$ in the Dimofte--Niu--Py
realization of the dg-shifted Yangian (cf.\ Vol~I,
Theorem~\ref{thm:dnp-bar-cobar-identification});
\item the prefactor $\hbar^{\mathrm{coll}}$ of the binary CY collision
residue $r_{CY}(z)$.
\end{enumerate}
The identification $\hbar^{\mathrm{KZ}} = \hbar^{\mathrm{DNP}} =
\hbar^{\mathrm{coll}}$ holds canonically once the CY direction is fixed,
and follows from the universality of the binary residue in
$\fg^{\mathrm{mod}}_{A_\cC}$.
\end{construction}

\begin{remark}[Convention sanity check]
\label{rem:three-hbar-sanity}
The identification of Construction~\ref{constr:three-parameter-hbar}
fails if any of the three parameters carries a hidden $2\pi i$ or
$1/n!$. In particular, the KZ parameter is conventionally normalized
so that the connection $\nabla = d - \hbar^{-1}\sum_i r_{ij}\,d\log(z_i
- z_j)$ has integer monodromy on the configuration space; the DNP
parameter is normalized so that the loop algebra has bracket
$[a_m, b_n] = [a, b]_{m+n} + m\hbar\delta_{m+n,0}\langle a, b\rangle$;
the collision residue is normalized so that $\Omega = r_{CY}(z)\cdot z$
at $z = 0$. These normalizations are mutually compatible; the
identification is parameter-by-parameter, not just up to scale.
\end{remark}

%% =====================================================================
%% SECTION 8: Face 7 -- Gaudin model from CY3
%% =====================================================================

\section{Face 7: Gaudin model from CY$_3$ collision residues}
\label{sec:face7-cy3-gaudin}

The seventh face is new in this volume. In Vol~I, face~7 of the standard
seven-face master is the Gaudin model arising from a chiral algebra
with primitive generator (Vol~I,
Theorem~\ref{thm:gaudin-yangian-identification}). In the CY setting, the
Gaudin model arises from CY$_3$ collision residues, with Gaudin
Hamiltonians constructed directly from the toric or elliptic $r$-matrix
of face~5 or face~6, and with the Bethe ansatz interpreted in terms of
DT/PT counts.

\begin{theorem}[Face 7 for $\C^3$: Gaudin from the Jordan-quiver residue]
\label{thm:face7-cy3-gaudin}
\ClaimStatusProvedHere
Let $X = \C^3$. Fix $n$ marked points $z_1, \ldots, z_n$ on the chiral
curve and Fock representations $V_1, \ldots, V_n$ of
$Y(\widehat{\fgl}_1)$. Define
\[
 H_i^{\C^3} \;:=\; \sum_{j \neq i} \frac{k\,\Omega_{ij}^{Y(\widehat{\fgl}_1)}}{z_i - z_j},
\]
acting on $V_1 \otimes \cdots \otimes V_n$, where
$\Omega_{ij}^{Y(\widehat{\fgl}_1)}$ is the Casimir tensor
$\Res_{z=0}\bigl[r_{CY}^{\C^3}(z)\bigr]$ acting on the $i$-th and
$j$-th factors. Then:
\begin{enumerate}[label=(\roman*)]
\item The operators $\{H_i^{\C^3}\}_{i=1}^n$ commute pairwise:
$[H_i^{\C^3}, H_j^{\C^3}] = 0$;
\item The joint spectrum is governed by the $\cW_{1+\infty}$ Bethe
ansatz of Litvinov and Maulik--Okounkov;
\item The DT counts of $\C^3$ (MacMahon plane-partition generating
function) provide the integer multiplicities of the Bethe roots: a
Bethe state at the configuration $\{w_a\}$ corresponds to a plane
partition of the appropriate shape.
\end{enumerate}
\end{theorem}

\begin{conjecture}[Face 7 for general toric/toroidal CY$_3$]
\label{conj:face7-general-gaudin}
The statement of Theorem~\ref{thm:face7-cy3-gaudin} extends to any
CY$_3$ $X$ admitting a toric (face~5) or toroidal (face~6) chiral
realization, with Hamiltonians $H_i^{X} := \sum_{j \neq i} r_{CY}^X(z_i
- z_j)_{ij}$, commutativity following from the classical Yang--Baxter
equation for $r_{CY}^X$, and DT/PT counts of $X$ supplying the Bethe
root multiplicities.
\ClaimStatusConjectured
\end{conjecture}

\begin{remark}[Proof sketch for $\C^3$]
For $X = \C^3$, $r_{CY}^{\C^3}(z) = k\,\Omega_{Y(\widehat{\fgl}_1)}/z$ has
a single pole. The Hamiltonians $H_i^{\C^3} = \sum_{j \neq i}
k\,\Omega_{ij}/(z_i - z_j)$ are the standard Gaudin Hamiltonians for
$Y(\widehat{\fgl}_1)$, and their commutativity follows from the
classical Yang--Baxter equation for $r_{CY}^{\C^3}$
(Construction~\ref{constr:c3-gaudin-from-residue}; cf.\ Vol~I,
Theorem~\ref{thm:gaudin-yangian-identification}). The Bethe ansatz
identification is the Litvinov $\cW_{1+\infty}$ Bethe ansatz, which
identifies Bethe roots with crystal melting configurations of $\C^3$
(equivalently, plane partitions).
\end{remark}

\begin{remark}[Why CY$_3$ supplies the Bethe roots]
\label{rem:cy3-bethe-content}
The novelty of face~7 in the CY setting is that the Bethe roots are not
abstract numbers; they are CY$_3$ DT/PT counts. For $\C^3$ the DT/PT
counts are MacMahon-type partition functions (3D plane partitions), and
each plane partition contributes a single Bethe root. For more general
toric CY$_3$, the topological vertex provides the partition counts and
hence the Bethe-root populations. In this sense the CY$_3$ Gaudin model
``geometrizes the Bethe ansatz'' in the same way that the toric quiver
geometrizes the affine Yangian.
\end{remark}

\begin{remark}[Face 7 is genuinely new for CY$_3$]
\label{rem:face7-genuinely-new}
Vol~I face~7 (Gaudin from chiral algebra with primitive generator) and
Vol~II face~7 (Gaudin from twisted holography boundary algebra) both
require an algebra with explicit generators. For CY$_3$, the analog
must extract Gaudin data from the conjectural quantum vertex chiral
group $G(X)$ via its CoHA presentation; the construction passes through
the $r$-matrix of face~5 (or face~6 in the toroidal case), not directly
through generators of $G(X)$ itself. This is the genuinely new content
of T4 in the CY setting and is the main object of this section.
\end{remark}

%% =====================================================================
%% SECTION 9: The seven-face master theorem
%% =====================================================================

\section{The seven-face master theorem in the CY setting}
\label{sec:cy-seven-face-master}

The preceding seven sections each realize the binary CY collision residue
in one geometric language. The master theorem asserts that all seven
realizations are canonically equivalent, with the equivalences given by
explicit specialization or limit, and with the limit/specialization data
recorded by the seven-face commuting diagram of Vol~I.

\begin{conjecture}[Seven faces of $r_{CY}(z)$]
\label{conj:cy-seven-face-master}
Let $\cC$ be a Calabi--Yau category of dimension $d \in \{2, 3\}$, and
let $A_\cC$ be its chiral algebraization (proved for $d = 2$ and $d = 3$; Theorem~CY-A$_3$, Theorem~\ref{thm:cy-to-chiral-d3}). Then the seven realizations
\[
 r_{CY}^{(\mathrm{F1})},\;
 r_{CY}^{(\mathrm{F2})},\;
 r_{CY}^{(\mathrm{F3})},\;
 r_{CY}^{(\mathrm{F4})},\;
 r_{CY}^{(\mathrm{F5})},\;
 r_{CY}^{(\mathrm{F6})},\;
 r_{CY}^{(\mathrm{F7})}
\]
constructed in
Sections~\ref{sec:face1-bar-cobar}--\ref{sec:face7-cy3-gaudin}
all coincide, in the sense that
\[
 r_{CY}^{(\mathrm{F}_i)} \;=\;
 \Phi_{i,j}\bigl(r_{CY}^{(\mathrm{F}_j)}\bigr)
 \qquad \forall\, i, j \in \{1, \ldots, 7\},
\]
for canonical change-of-language morphisms $\Phi_{i,j}$ given explicitly
by Theorem~CY-A$_2$ in the proved $d = 2$ case, and conjecturally by
CY-A$_3$ (face~1 to face~3), the
Schiffmann--Vasserot/RSYZ identification (face~2 to face~5), the MO
stable envelope (face~4), the elliptic specialization $\hbar \to 0$
(face~6 to face~3), and the Litvinov Bethe ansatz (face~5 to face~7).
The limits between faces are recorded in the seven-face commuting
diagram (cf.\ Vol~I,
Theorem~\ref{thm:vol1-seven-face-master}; Vol~II,
Theorem~\ref{thm:vol2-seven-face-master}).
\end{conjecture}

\begin{remark}[Status by face]
\begin{description}
\item[Face 1] (CY bar-cobar twisting): \ClaimStatusProvedHere{} for
$d = 2$ and $d = 3$ (CY-A$_3$ proved, Theorem~\ref{thm:cy-to-chiral-d3}).
\item[Face 2] (CoHA): \ClaimStatusProvedElsewhere{} for the existence
of the CoHA (Kontsevich--Soibelman, Davison--Meinhardt);
identification with $r_{CY}$ is \ClaimStatusConjectured{} in general,
\ClaimStatusProvedHere{} for $\C^3$ (Schiffmann--Vasserot,
Theorem~\ref{thm:sv-c3}) and toric without compact $4$-cycles (RSYZ,
Theorem~\ref{thm:rsyz}).
\item[Face 3] (classical coisson, genus~$0$ only): \ClaimStatusProvedHere{} for $d = 2$
via the proved coisson--PVA dictionary; \ClaimStatusConjectured{} for
$d = 3$.
\item[Face 4] (MO $R$-matrix): \ClaimStatusProvedElsewhere{} for the
existence of $R_{\mathrm{MO}}$ (Maulik--Okounkov 2019);
identification with $r_{CY}$ for $K3 \times E$ is~\ClaimStatusConjectured{}
in general, with explicit verification at the Gepner self-dual limit
(Conjecture~\ref{conj:k3e-self-dual-limit}).
\item[Face 5] (toric Yangian): \ClaimStatusProvedHere{} for $\C^3$;
\ClaimStatusProvedElsewhere{} for toric without compact $4$-cycles via
RSYZ.
\item[Face 6] (elliptic Sklyanin): \ClaimStatusProvedElsewhere{} for the
classical Sklyanin bracket and its quantization to Felder's elliptic
quantum group; identification with $r_{CY}$ is~\ClaimStatusProvedHere{}
in the $\fsl_2$ Heisenberg case via the Fay identity
(Proposition~\ref{prop:fay-implies-d-squared}).
\item[Face 7] (CY$_3$ Gaudin): \ClaimStatusProvedHere{} for $\C^3$
(Construction~\ref{constr:c3-gaudin-from-residue}); \ClaimStatusConjectured{}
for general toric/toroidal CY$_3$.
\end{description}
The status partition above is consistent: each face's
classical components are tagged ProvedElsewhere with explicit
attribution, and only the genuinely new identifications carry the
ProvedHere tag.
\end{remark}

\begin{remark}[Three invariants across the seven faces]
\label{rem:three-invariants-cross-faces}
By Remark~\ref{rem:cy-three-invariants}, every realization of $r_{CY}$
carries the three invariants $(p_{\max}, k_{\max}, r_{\max})$.
AP59 requires that each face state which invariant is being measured, and
forbids identifying them.
\begin{itemize}
\item Face~5 (toric Yangian) is sensitive to $p_{\max}$: the pole
structure of $r_{Y(\widehat{\fg})}$ is the OPE pole structure of the
toric quiver superpotential.
\item Face~6 (elliptic) carries $k_{\max}$ in its theta-function
expansion: at $k_{\max} = 0$ the residue is purely a single-pole
elliptic term; at $k_{\max} \geq 1$ there are derivative terms.
\item Face~1 (bar-cobar) is sensitive to $r_{\max}$: the shadow tower
$\Theta_A^{\leq r}$ terminates at degree $r_{\max}$, and the binary
residue is the projection to $r = 2$.
\end{itemize}
The seven-face master statement is that the binary residue is the
\emph{same} object across all faces, but the sensitivity to the three
invariants differs face by face.
\end{remark}

%% =====================================================================
%% SECTION 10: Cross-volume bridge
%% =====================================================================

\section{Cross-volume bridge}
\label{sec:cy-cross-volume-bridge}

Part~\ref{part:connections} of Vol~III mirrors the seven-face programme
developed in Vols.~I and II. The common skeleton is the seven-face
programme; the variation is in which face is most concrete.

\begin{remark}[The three seven-face masters]
\label{rem:three-seven-face-masters}
The three volumes each devote a part to the seven-face programme, with
the same architecture but different ground objects:
\begin{enumerate}[label=(\arabic*)]
\item \emph{Vol~I, the seven-face part:} the binary collision residue of a chiral
algebra on a curve, in seven languages: bar-cobar twisting, primitive
generator, classical $r$-matrix, KZ connection, Gaudin model, Bethe
ansatz, dg-shifted Yangian (cf.\ Vol~I,
Theorem~\ref{thm:vol1-seven-face-master}).
\item \emph{Vol~II, the seven-face part:} the binary collision residue of a
holomorphic-topological quantum group, in seven languages: open-string
brace algebra, derived center, twisted holography boundary, line
defect, Wilson line $R$-matrix, soft graviton coupling, Yangian double
(cf.\ Vol~II, Theorem~\ref{thm:vol2-seven-face-master}).
\item \emph{Vol~III, this chapter:} the binary CY collision residue of a
Calabi--Yau chiral algebra, in seven CY-specific languages: CY
bar-cobar, CoHA / perverse coherent sheaves, classical CY Poisson
coisson, MO stable envelope, affine super Yangian for toric CY$_3$,
elliptic Sklyanin for toroidal CY, Gaudin from CY$_3$
(Conjecture~\ref{conj:cy-seven-face-master} above).
\end{enumerate}
The three master theorems are mutually compatible: under the CY-to-chiral
functor $\Phi$, face~$i$ of Vol~III maps to a specialization of face~$i$
of Vol~I, and similarly for Vol~II. The CY setting is the most
constrained: each face acquires geometric content that the abstract
chiral algebra setting does not see (DT counts, plane partitions,
crystal melting, MO stable envelopes).
\end{remark}

\begin{remark}[The CY direction is the new slot]
\label{rem:cy-direction-new-slot}
Compared with Vol~I and Vol~II, the new feature of the CY seven-face
master is the presence of the CY direction $\Omega_\cC$. In Vol~I the
binary residue lives on $\overline{C}_2(X) \times \fg^{\otimes 2}$; in
Vol~II it lives on $\overline{C}_2(X) \times \mathrm{open\text{-}sector}^{\otimes 2}$;
in Vol~III it lives on $\overline{C}_2(X) \times \cC^{\otimes 2}$, with
the CY volume form supplying the pairing on $\cC^{\otimes 2}$. The CY
direction is what makes face~4 (MO) and face~7 (CY$_3$ Gaudin) genuinely
geometric: they live on quiver varieties whose definition requires the
CY structure.
\end{remark}

\begin{remark}[$\kappa_{\mathrm{ch}}(\cA)$ is a face-by-face invariant]
\label{rem:kappa-face-by-face}
The modular characteristic $\kappa_{\mathrm{ch}}(A_\cC)$ depends on the construction,
not on the geometry alone (AP-CY59). For $K3 \times E$, the $\kappa_\bullet$-spectrum
has four values from four distinct sources (AP-CY68: total-space vs fiber):
$\kappa_{\mathrm{cat}} = \chi(\cO_{K3 \times E}) = 2 \cdot 0 = 0$ (K\"unneth;
the K3-fiber value $\chi(\cO_{K3}) = 2$ is a distinct \emph{fiber} invariant,
not $\kappa_{\mathrm{cat}}$ of the total space) and $\kappa_{\mathrm{fiber}} = 24$ are
\emph{manifold invariants} (topological, independent of any construction),
while $\kappa_{\mathrm{ch}} = 3$ (from the CY-to-chiral functor~$\Phi$) and
$\kappa_{\mathrm{BKM}} = 5$ (from the Borcherds lift) are
\emph{construction-dependent invariants} arising from distinct mathematical
constructions, not from four applications of~$\Phi$ (preface).
Each value sits on a different
face: $\kappa_{\mathrm{ch}}$ on face~1 (bar-cobar of the chiral de Rham
complex), $\kappa_{\mathrm{cat}}$ on face~2 (holomorphic Euler characteristic),
$\kappa_{\mathrm{fiber}}$ on face~3 (classical lattice VOA limit),
$\kappa_{\mathrm{BKM}}$ on face~4 (MO/Igusa). The $\kappa_{\mathrm{ch}}$-spectrum is
the imprint of the geometry on the seven-face master. The rule is that $\kappa_{\mathrm{ch}}(\cA)$ depends on the full algebra, not on its Virasoro
subalgebra; in the CY setting, $\kappa_{\mathrm{ch}}(\cA)$ depends on the full
\emph{construction}, not on the underlying CY manifold.
\end{remark}

\begin{remark}[Outlook: an eighth face?]
\label{rem:outlook-eighth-face}
The standard seven-face master is closed: there are exactly seven
languages in which the binary residue lives, corresponding to the seven
edges of the Vol~I commuting diagram. The natural question in the CY setting: does the
deformation-quantization face (Khan--Zeng 3d PVA sigma model)
constitute an eighth? At present this face reduces to a
combination of face~3 (classical coisson) and face~7 (Gaudin from
CY$_3$), with the deformation-quantization parameter playing the role of
the loop parameter in face~7. An independent eighth face would require
a CY-specific structure not visible in the seven-face commuting diagram;
it remains open whether the
$3$d~$\cN = 2$ holomorphic-twisted theory provides such a structure.
\end{remark}

%% =====================================================================
%% Closing
%% =====================================================================

The genus-$1$ extension, identifying the KZB connection with the
elliptic $r$-matrix on the torus~$E_\tau$, is proved for affine KM
in Vol~I, Chapter~\texttt{ch:genus1-seven-faces}.

\noindent\textit{This chapter closes the seven-face programme for
Vol~III. Subsequent chapters in Part~\ref{part:frontier} record the geometric Langlands
implications and the bridges to Vols~I--II at the level of theorem
statements; the present chapter is the algebraic engine that makes the
bridges possible.}
